Main object to apply kNN nowadays are embeddings (vectors) produced by neural networks.
Such embeddings represent object properties in a compact format.
Nearest neighbors are used to roughly select possible candidates for further processing.

\subsection{Зачем нужны индексы}

KNN index is a data structure that allows you to quickly find the $N$ nearest vectors for a given vector.
We perform some precompute beforehand to increase speed of final search.

\subsection{Идеи индексов для kNN от простых к сложным}

\begin{enumerate}
    \item Clustering
    \item k-d tree
    \begin{itemize}
        \item default for solid team
        \item ball tree is extension for high dimensional cases
    \end{itemize}
    \item HNSW (most common nowadays)
\end{enumerate}

\subsection{k-dimensional tree (kd-tree)}

Представим на секунду, что у нас есть всего лишь один признак, то есть объекты выражаются вещественными числами, а
не векторами.
В этом случае для поиска ближайшего соседа напрашивается всем вам известное бинарное дерево поиска, которое
позволяет находить элементы за логарифмическое время.
Оказывается, существует аналог данной структуры в многомерном пространстве, который называется k-d-дерево
(k-d tree, сокращение от k-dimensional tree).

Как и в обычном дереве поиска, в k-d-дереве каждый узел является объектом обучающей выборки, который особым образом
делит пространство на два полупространства.
Таким образом, всё пространство оказывается поделено на множество малых областей, и такое деление оказывается очень
полезным при поиске ближайших соседей.

Рассмотрим подробнее, как строится такое дерево.
Трудность в применении обычного дерева поиска состоит в том, что мы не можем напрямую сравнить два вектора так же,
как два вещественных числа.
Чтобы эту проблему преодолеть, узлы дерева будут делить пространство лишь по одной оси.
При движении вниз по дереву оси, по которым точки делят пространство, циклически сменяют друг друга.
Например, в двумерном пространстве корень будет отвечать за деление по $x$-координате, его сыновья -- за деление
по y-координате, а внуки -- снова за x-координату, и так далее.
Посмотрим, как это работает на примере:

На картинке выше корень $(30, 40)$ делит все точки по оси $x$: слева оказываются точки, у которых $x < 30$,
а справа -- те, у которых $x \geq 30$.
Аналогично левый сын корня $(5,25)$ делит своё поддерево по оси $y$: слева оказываются точки, у которых $y<25$,
а справа -- те, у которых $y \geq 25$.

Остаётся вопрос -- как выбирать точки, которые будут делить пространство пополам?
Чтобы дерево было сбалансированным, нужно находить точку с медианой, соответствующей уровню поддерева координаты.
На практике часто ограничиваются выбором случайной точки или любой эвристикой по приближённому поиску медианы
(например, медиана некоторого подмножества точек).
Это позволяет ускорить построение дерева, но убирает все гарантии на его сбалансированность.

Добавлять новые точки можно так же, как и в одномерном дереве поиска.
Спускаясь по дереву, можно однозначно определить лист, к которому нужно подвесить новую точку, чтобы не нарушить все
свойства дерева.
При добавлении большого количества точек, однако, дерево может перестать быть сбалансированным, и нужно проводить
ребалансировку.
Также существуют варианты k-d-деревьев, которые сохраняют сбалансированность при добавлении / удалении точек.

Поговорим теперь про то, как же находить ближайших соседей с помощью такого дерева.
Будем производить обход дерева в глубину с двумя модификациями.

\begin{itemize}
    \item Во-первых, будем запоминать наиболее близкую точку.
    Это позволит не заходить в поддеревья, задающие области, которые заведомо дальше, чем текущая наиболее близкая
    точка, поэтому не имеет смысла искать в них ближайших соседей.
    \item Во-вторых, будем прежде всего обходить те поддеревья, которые задают наиболее близкие области, а значит, с
    большей вероятностью содержат ближайшего соседа.
\end{itemize}

Сложность метода по размеру обучающей выборки в среднем равна $O(\log N)$ при равномерном распределении точек.
При большой размерности пространства, однако, алгоритму приходится посещать больше ветвей дерева, чтобы найти
ближайших соседей.
Например, если $N \approx D$, то сложность становится примерно такой же, как и в случае полного перебора.
В общем случае для того чтобы асимптотика действительно была логарифмической, нужно, чтобы $N \geq 2^D$.
Поэтому уже при количестве признаков порядка сотни алгоритм не даёт существенных преимуществ перед полным перебором

\subsection{Annoy}

Алгоритмы, основанные на деревьях, очень часто применяются в задачах поиска соседей.
Идея всех таких методов заключается в итеративном разделении пространства случайными гиперплоскостями и построении
на базе этого разделения дерева, в листах которого содержится малое число объектов.

Одним из наиболее ярких представителей этого семейства является Annoy -- алгоритм, который используется Spotify
для рекомендаций музыки.
Задача подобных рекомендательных систем довольно простая -- нужно посоветовать пользователю композиции, которые
он ещё не слушал, но которые при этом с высокой вероятностью ему понравятся.
Простая и рабочая идея -- предлагать композиции, похожие на те, которые он уже слушает.
Здесь на помощь как раз и приходят методы поиска ближайших соседей.

Annoy в какой-то степени похож на k-d-деревья.
Сначала выбираются два случайных объекта обучающей выборки и проводится гиперплоскость, симметрично их разделяющая.
Затем для каждого полученного полупространства итеративно запускается такая же процедура, которая продолжается до
тех пор, пока в каждой области будет не более $M$ объектов ($M$ — гиперпараметр).

Таким образом задаётся бинарное дерево с глубиной порядка $O(logN)$ в среднем.
Спускаясь по этому дереву, можно найти область, в которой лежит целевой объект и некоторое количество близких к нему
элементов обучающей выборки.
Проблема в том, что это не обязательно будут самые близкие объекты, поэтому для увеличения точности составляется лес
из таких деревьев и берётся объединение соответствующих целевому объекту областей.

Чем больше таких деревьев берётся, тем более точным будет результат, но придётся тратить большее время на его поиск.

Преимущество алгоритма -- простота нахождения компромисса между скоростью работы и точностью с помощью тюнинга
гиперпараметров.
К минусам можно отнести то, что алгоритм плохо параллелится и переносится на GPU, не работает эффективно с батчами,
а также то, что для добавления новой точки в обучающую выборку придётся перезапускать процедуру с самого начала

\subsection{HNSW}

Следующий класс алгоритмов основан на построении специального графа близости (proximity graph) на объектах выборки
и дальнейшем жадном поиске по этому графу.
Алгоритмы этого семейства сейчас считаются state-of-the-art (SotA) для многих задач.

Рассмотрим подробнее этот класс алгоритмов на примере одного из наиболее популярных из них под названием
Navigable small world (NSW).
Идея его в следующем: на данных строится граф (он также называется NSW), который удовлетворяет двум следующим
свойствам:

\begin{enumerate}
    \item Между любыми двумя точками существует короткий путь, или, более формально, матожидание числа кратчайшего
    пути между двумя случайно выбранными вершинами растёт как $O(logN)$.
    \item Средняя степень вершины мала.
\end{enumerate}

На первый взгляд может показаться, что тяжело выполнить одновременно оба свойства, но на самом деле большая часть
графов в реальном мире являются NSW-графами.
Самый простой пример -- это известное правило шести рукопожатий: любые два случайных человека соединены короткой
последовательностью личных контактов длиной не более шести, несмотря на то, что количество знакомых у среднего
человека (100-1000) мало по сравнению с населением Земли.

В таких графах существует очень простой метод поиска соседей.
Нужно выбрать случайную точку, среди её соседей выбрать того, который ближе всего к целевому объекту, и повторить
процедуру уже для него.
Показано, что такой жадный поиск имеет полилогарифмическую асимптотику.

\subsubsection{Построение графа}

Проблема такого подхода в том, что можно попасть в плотный кластер и очень долго оттуда выбираться.
Для решения этой проблемы используется иерархия NSW, или \textit{Hierarchical navigable small world (HNSW)}.
Исходный граф является нулевым слоем.
Каждый следующий слой строится в два шага

\begin{enumerate}
    \item Каждая вершина текущего слоя попадает в следующий с некоторой вероятностью $p$.
    \item На всех вершинах, попавших в новый слой, строится NSW\@.
\end{enumerate}
По построению количество слоёв будет $O(logN)$.

Поиск начинается в самом верхнем слое.
После нахождения ближайшей к целевому объекту вершины спускаемся на слой ниже и начинаем поиск из этой вершины.
Повторяем процедуру, пока не спустимся до нулевого слоя.
Таким образом, на каждом слое мы всё больше уточняем наш ответ.
Стоит отметить, что для ускорения работы иногда поиск останавливают не при нахождении ближайшей вершины, а раньше,
используя критерии остановки.

Интуитивно легко понять, почему такая иерархическая структура решает проблему плотных кластеров: в верхних слоях
вершин мало, а расстояния между ними в среднем большие, а значит, таких кластеров там почти нет.
Поэтому, попадая в нижний слой, мы чаще всего оказываемся уже в нужном кластере и просто уточняем результат работы
алгоритма.

\subsubsection{Использование графа}

HNSW, так же как и рассмотренные ранее приближённые методы, позволяет искать трейд-офф между точностью и скоростью.
Плюс ко всему на реальных данных он часто работает лучше других методов и сейчас считается предметом искусства.
Однако этот способ поиска не лишён и недостатков.
Главный заключается в том, что нельзя добавлять точки в обучающую выборку без перестройки структуры.
Помимо этого, он довольно требователен по памяти из-за того, что для каждого слоя приходится хранить как вершины,
которые в него входят, так и связи между этими вершинами

\subsection{Примеры векторных баз данных}

\begin{itemize}
    \item OpenAI
    \item Qdrant
    \item Try chroma
    \item Feast
    \item NVIDIA AI store
\end{itemize}
Все, кроме первой, ready for production